% ==============================================================================
% METADATOS DE AUDITORÍA EDITORIAL Y RIGOR CIENTÍFICO (PROTOCOLO AGENTS.md)
% ==============================================================================
% Capítulo: 01
% Título: Introducción a la Inteligencia Artificial: IA Generativa
% Grupo de Origen: Grupo 11 (Asignatura de Inteligencia Artificial Aplicada, UNMSM)
% Archivo Fuente Analizado: Capitulos IA S2/Grupo 11/Capitulo1_IntroducciónIAG.pptx (24 diapositivas)
% Enfoque: Inteligencia Artificial Aplicada (traslado fiel letra por letra e imagen por imagen)
% Cumplimiento Normativo: APA 7ma Edición + House Style UNMSM (7 componentes pedagógicos canónicos).
% ==============================================================================

\section[Introducción a la Inteligencia Artificial. IA Generativa]{Introducción a la Inteligencia Artificial. Inteligencia Artificial Generativa}
\markboth{CAPÍTULO 1\quad INTRODUCCIÓN A LA IA: IA GENERATIVA}{}

\begin{cajabeige}
\textbf{Pregunta Central del Capítulo:} \textquestiondown Cómo evolucionaron los modelos y las arquitecturas de la IA generativa (Transformers, redes generativas antagónicas y modelos de difusión) y qué capacidades y límites plantean para comprenderla y usarla en la práctica?\\[0.25cm]
\textbf{Objetivo de Aprendizaje:} En este capítulo vas a descubrir qué es la IA y sus tipos, cómo se llegó a los sistemas generativos actuales, cómo aprende una máquina a partir de datos y por qué la IA generativa marca un antes y un después frente a la IA tradicional.
\end{cajabeige}

\textoside{
Durante décadas, las computadoras siguieron reglas fijas y ordenaron información que ya existía. Hoy, la inteligencia artificial generativa (GenAI) cambió el juego. Basta con escribir una orden en cualquier idioma para que un sistema cree un texto, un fragmento de código o una imagen que antes no existía. Eso abre preguntas interesantes: \textquestiondown puede una máquina ser creativa? \textquestiondown Hasta dónde llega esta tecnología? \textquestiondown Cómo conviene usarla para estudiar y trabajar sin perder el juicio crítico?
}{CONCEPTO CLAVE\\[0.1cm]IA Generativa}

\subsubsection*{Qué encontrarás en el capítulo}
\begin{itemize}
    \item \textbf{\textquestiondown Qué es la Inteligencia Artificial?:} Definición, usos cotidianos, IA estrecha y general.
    \item \textbf{Breve historia de la IA:} Línea de tiempo, personajes e hitos por etapas.
    \item \textbf{Cómo aprende una máquina:} Aprendizaje automático, redes neuronales, datos y entrenamiento.
    \item \textbf{La IA generativa:} De clasificar a generar, modelos de lenguaje y un caso de aplicación.
\end{itemize}

\vspace{0.3cm}

\subsection{Qué es la Inteligencia Artificial}

Es probable que hoy hayas desbloqueado tu teléfono con el rostro, recibido una sugerencia de video o usado un traductor automático sin detenerte a pensar que detrás hay inteligencia artificial. La expresión se volvió común, pero también ambigua: se usa para hablar de un filtro de correo, de un robot de película y de un asistente que redacta ensayos. Antes de estudiar la IA generativa conviene fijar con claridad de qué hablamos.

\subsubsection{Una definición sencilla}
La inteligencia artificial (IA) es el campo de la informática que diseña sistemas capaces de realizar tareas que, si las hiciera una persona, diríamos que requieren inteligencia: percibir su entorno, aprender de la experiencia, razonar, usar el lenguaje y tomar decisiones. 

La idea es antigua. La propuesta que dio origen al taller de Dartmouth sostenía que el aprendizaje, el uso del lenguaje, la resolución de problemas y la creatividad podían describirse con tanta precisión que una máquina podría simularlos. Años antes, Turing (1950) había reemplazado la pregunta "\textquestiondown pueden pensar las máquinas?" por una prueba observable: el "juego de imitación".

Importa notar que el término no tiene una definición única y universal. García-Peñalvo y Vázquez-Ingelmo (2023) lo constatan incluso para la IA generativa, tras revisar 631 soluciones publicadas en cinco años. Por eso Feuerriegel et al. (2024) proponen mirarla en tres niveles:
\begin{enumerate}
    \item \textbf{Los modelos:} Los algoritmos entrenados.
    \item \textbf{Los sistemas:} Que los empaquetan con datos e interfaces.
    \item \textbf{Las aplicaciones:} Los usos concretos en cada sector.
\end{enumerate}

\begin{cajagris}
\textbf{CONCEPTO CLAVE: Inteligencia Artificial}\\
Campo que busca construir sistemas capaces de realizar tareas que normalmente asociamos con la inteligencia humana.
\end{cajagris}

\subsubsection{Qué tipos de tareas aborda la IA}
Más que una sola técnica, la IA es un conjunto de métodos que se aplican a distintas familias de tareas, resumidas en la Tabla 1.1.

\begin{table}[htpb!]
\centering
\small\sffamily
\caption{\textbf{Tabla 1.1.} \textit{Familias de tareas de la IA}}
\vspace{0.15cm}
\begin{tabularx}{\textwidth}{>{\bfseries}p{3.2cm} >{\raggedright\arraybackslash}p{4.2cm} >{\raggedright\arraybackslash}p{4.2cm} >{\raggedright\arraybackslash}X}
\toprule
Tarea & Qué hace el sistema & Ejemplo & Datos de entrada \\
\midrule
Percibir & Interpreta imágenes, sonidos u otras señales. & Reconocimiento facial, dictado por voz. & Imágenes, audio \\
\midrule
Comprender lenguaje & Analiza, clasifica o traduce texto. & Traductor automático, filtro de correo. & Texto \\
\midrule
Predecir y decidir & Estima un valor o elige entre opciones. & Recomendaciones, detección de fraude. & Tablas, historiales \\
\midrule
Generar & Crea contenido nuevo a partir de lo aprendido. & Asistentes de redacción, imágenes sintéticas. & Texto, imagen, audio, código \\
\bottomrule
\end{tabularx}
\vspace{0.15cm}
\raggedright\scriptsize \textbf{Fuente:} Elaboración propia a partir del material del Grupo 11.
\end{table}

\subsubsection{Lo que la IA no es}
Gran parte de los malentendidos nace de atribuirle a la IA cualidades que no tiene. Cinco ideas ayudan a corregirlos:
\begin{itemize}
    \item \textbf{No es una mente consciente:} Un modelo de lenguaje produce texto coherente, pero eso no significa que comprenda lo que dice.
    \item \textbf{No es un buscador ni una base de datos:} Un sistema generativo no "recupera" respuestas guardadas: las construye a partir de patrones aprendidos.
    \item \textbf{No es infalible:} Puede producir información falsa con aparente seguridad, fenómeno conocido como alucinación.
    \item \textbf{No es neutral:} Hereda los sesgos de los datos con los que se entrenó y puede reproducirlos o amplificarlos.
    \item \textbf{No es una caja mágica:} Entender cómo funciona permite evaluarla críticamente; por eso se habla de "descajanegrización".
\end{itemize}

\subsubsection{Un ejemplo sencillo: el filtro de correo no deseado}
El filtro de spam de tu correo es un buen ejemplo de IA que no genera nada, pero aprende. Su funcionamiento se puede resumir en cuatro pasos:
\begin{enumerate}
    \item Se reúnen miles de correos que personas ya clasificaron como "spam" o "no spam".
    \item El sistema busca patrones: ciertas palabras, enlaces sospechosos, remitentes poco habituales.
    \item Con lo aprendido, asigna a cada correo nuevo una probabilidad de ser spam.
    \item Si la probabilidad supera un umbral, el correo va a la carpeta de spam; las correcciones del usuario sirven para mejorar el sistema.
\end{enumerate}

Este sistema clasifica. Un asistente generativo haría algo distinto: redactaría un correo nuevo. Esa diferencia se retoma en la sección 1.4.

\begin{figure}[htpb!]
\centering
\begin{tcolorbox}[colback=white, colframe=unmsm_gold, arc=3pt, boxrule=0.8pt, top=0.3cm, bottom=0.3cm, left=0.4cm, right=0.4cm]
\centering
\small\sffamily
\textbf{Figura 1.1.} \textit{La IA y sus subcampos: de lo general a lo específico}\\[0.25cm]
\resizebox{0.92\linewidth}{!}{%
\begin{tikzpicture}[>=stealth, font=\small\sffamily]
    % Caja IA
    \draw[rounded corners=6pt, fill=bg_marfil, draw=unmsm_red, line width=1.2pt] (-6.5,-5.2) rectangle (6.5,5.2);
    \node[anchor=north west, text=unmsm_red, font=\normalsize\bfseries] at (-6.1,4.8) {Inteligencia Artificial (IA)};
    \node[anchor=north west, text=unmsm_gray, text width=12cm] at (-6.1,4.2) {Sistemas que realizan tareas que asociamos con la inteligencia humana.};

    % Caja ML
    \draw[rounded corners=5pt, fill=bg_arena, draw=unmsm_darkred, line width=1.0pt] (-5.8,-4.6) rectangle (5.8,3.2);
    \node[anchor=north west, text=unmsm_darkred, font=\small\bfseries] at (-5.4,2.8) {Aprendizaje automático (ML)};
    \node[anchor=north west, text=black, text width=10.8cm] at (-5.4,2.2) {Sistemas que aprenden patrones a partir de datos, sin reglas escritas a mano.};

    % Caja DL
    \draw[rounded corners=4pt, fill=white, draw=unmsm_red, line width=0.8pt] (-5.1,-4.0) rectangle (5.1,0.9);
    \node[anchor=north west, text=unmsm_red, font=\small\bfseries] at (-4.7,0.5) {Aprendizaje profundo (DL)};
    \node[anchor=north west, text=black, text width=9.5cm] at (-4.7,-0.1) {Aprendizaje con redes neuronales de muchas capas.};

    % Caja GenAI
    \draw[rounded corners=3pt, fill=bg_marfil, draw=unmsm_darkred, line width=0.8pt] (-4.4,-3.4) rectangle (4.4,-1.0);
    \node[anchor=north west, text=unmsm_darkred, font=\small\bfseries] at (-4.0,-1.4) {IA generativa (GenAI)};
    \node[anchor=north west, text=black, text width=8.2cm] at (-4.0,-2.0) {Modelos que producen contenido nuevo: texto, imagen, audio o código.};
\end{tikzpicture}%
}\\[0.2cm]
\raggedright\scriptsize \textbf{Fuente:} Elaboración propia a partir de LeCun, Bengio y Hinton (2015) y Feuerriegel et al. (2024).
\end{tcolorbox}
\end{figure}

\subsubsection{La IA en la vida cotidiana}
La IA ya forma parte de rutinas que casi no notamos. Algunos ejemplos habituales:
\begin{itemize}
    \item \textbf{En el teléfono:} Desbloqueo facial, teclado predictivo, dictado por voz y traducción instantánea.
    \item \textbf{En plataformas digitales:} Recomendaciones de videos, música y productos según lo que ya consumiste.
    \item \textbf{En el correo y la banca:} Filtros de correo no deseado y detección de operaciones sospechosas.
    \item \textbf{En el transporte:} Rutas que se recalculan según el tráfico estimado.
    \item \textbf{En la salud:} Apoyo al análisis de imágenes médicas, siempre bajo supervisión de especialistas.
    \item \textbf{En la educación:} Asistentes conversacionales que explican, resumen o proponen ejercicios.
\end{itemize}

\subsubsection{IA estrecha frente a IA general}
Se suele distinguir dos grandes niveles de ambición, resumidos en la Tabla 1.2.

\begin{table}[htpb!]
\centering
\small\sffamily
\caption{\textbf{Tabla 1.2.} \textit{Comparación entre IA estrecha e IA general}}
\vspace{0.15cm}
\begin{tabularx}{\textwidth}{>{\bfseries}p{2.8cm} >{\raggedright\arraybackslash}p{6.0cm} >{\raggedright\arraybackslash}X}
\toprule
Criterio & IA estrecha (ANI) & IA general (AGI) \\
\midrule
Alcance & Resuelve tareas concretas para las que fue entrenada. & Resolvería cualquier tarea intelectual que pueda hacer una persona. \\
\midrule
Ejemplos & Reconocer rostros, traducir, recomendar, jugar ajedrez, redactar textos. & Todavía no existe; es un objetivo de investigación y objeto de debate. \\
\midrule
Transferencia & Rinde mal fuera de su dominio sin reentrenamiento. & Aprendería una tarea nueva y aplicaría lo aprendido a otras. \\
\midrule
Estado actual & Es toda la IA que se usa hoy. & Hipotética; no hay consenso sobre cómo ni cuándo se lograría. \\
\bottomrule
\end{tabularx}
\vspace{0.15cm}
\raggedright\scriptsize \textbf{Fuente:} Elaboración propia a partir del material del Grupo 11.
\end{table}

Los grandes modelos de lenguaje desdibujan un poco esta frontera, porque un mismo modelo puede traducir, resumir, programar y responder preguntas. Bommasani et al. (2021) destacan justamente que estos modelos se reutilizan en muchas aplicaciones. Sin embargo, esa versatilidad no equivale a inteligencia general: siguen dependiendo de los datos de entrenamiento y fallan de maneras que una persona no esperaría.

\subsubsection{Vocabulario necesario}
Estos términos aparecerán a lo largo del libro. Conviene tenerlos presentes desde ahora (Tabla 1.3).

\begin{table}[htpb!]
\centering
\small\sffamily
\caption{\textbf{Tabla 1.3.} \textit{Glosario básico de términos en Inteligencia Artificial}}
\vspace{0.15cm}
\begin{tabularx}{\textwidth}{>{\bfseries}p{2.8cm} >{\raggedright\arraybackslash}X}
\toprule
Término & Significado \\
\midrule
Dato & Cada ejemplo con el que se entrena o se usa un sistema: un texto, una imagen, una fila de una tabla. \\
\midrule
Algoritmo & Procedimiento paso a paso para resolver un problema o ajustar un modelo. \\
\midrule
Modelo & Función con parámetros ajustados durante el entrenamiento, que transforma entradas en salidas. \\
\midrule
Entrenamiento & Proceso de ajustar los parámetros del modelo para reducir su error sobre los datos. \\
\midrule
Inferencia & Uso del modelo ya entrenado sobre datos nuevos. \\
\midrule
Token & Pieza mínima de texto que procesa un modelo de lenguaje: una palabra o un fragmento. \\
\midrule
Prompt & Instrucción o texto de entrada que se le da a un modelo generativo. \\
\bottomrule
\end{tabularx}
\vspace{0.15cm}
\raggedright\scriptsize \textbf{Fuente:} Elaboración propia.
\end{table}

\begin{cajabeige}
\textbf{ACTIVIDAD PARA PENSAR:}\\
\textquestiondown Qué tareas de tu día a día resuelve hoy una IA estrecha? \textquestiondown Cuál de ellas te sorprendería que dejara de funcionar, y por qué?
\end{cajabeige}

\vspace{0.3cm}

\subsection{Breve historia de la IA}

La IA no nació con ChatGPT. Su recorrido alterna épocas de entusiasmo con periodos de desilusión, los llamados "inviernos de la IA", y su desarrollo reciente en la IA generativa solo se entiende como resultado de esa historia. La Figura 1.2 resume los hitos en una línea de tiempo; en las páginas siguientes se explican por etapas.

\subsubsection{Cómo leer esta sección}
Dividiremos la historia en cinco etapas. Para cada una responderemos tres preguntas: \textquestiondown qué ideas la hicieron posible?, \textquestiondown quiénes fueron sus protagonistas? y \textquestiondown qué dejó para la etapa siguiente? La Tabla 1.4 ofrece un resumen que puedes consultar mientras avanzas.

\begin{table}[htpb!]
\centering
\small\sffamily
\caption{\textbf{Tabla 1.4.} \textit{Resumen de personajes e hitos por etapa histórica}}
\vspace{0.15cm}
\begin{tabularx}{\textwidth}{>{\bfseries}p{4.0cm} >{\raggedright\arraybackslash}p{4.5cm} >{\raggedright\arraybackslash}X}
\toprule
Etapa & Personajes & Hito principal \\
\midrule
1. Fundamentos (1943--1955) & McCulloch, Pitts, Turing & Neurona artificial; test de Turing. \\
\midrule
2. Nacimiento (1956--1973) & McCarthy, Minsky, Rochester, Shannon, Rosenblatt & Taller de Dartmouth; perceptrón. \\
\midrule
3. Inviernos (1974--1993) & Rumelhart, Hinton, Williams & Sistemas expertos; retropropagación. \\
\midrule
4. Aprendizaje profundo (1997--2016) & LeCun, Bengio, Hinton, Goodfellow & Redes profundas; GAN. \\
\midrule
5. Era generativa (2017--hoy) & Vaswani, Brown, Ho, Ouyang & Transformer; GPT-3; difusión; RLHF. \\
\bottomrule
\end{tabularx}
\vspace{0.15cm}
\raggedright\scriptsize \textbf{Fuente:} Elaboración propia a partir de Nasser-Eddine (2024).
\end{table}

\begin{cajagris}
\textbf{LECCIÓN HISTÓRICA: Ciclos de entusiasmo y desilusión}\\
Cada avance genera expectativas altas. Los inviernos ocurrieron cuando esas expectativas superaron lo que los métodos podían ofrecer; conviene recordarlo al evaluar la IA generativa actual.
\end{cajagris}

\subsubsection*{Preguntas guía}
\begin{itemize}
    \item \textquestiondown Qué diferencia hay entre "simular" la inteligencia y "tenerla"? \textquestiondown Cómo responde el test de Turing?
    \item \textquestiondown Por qué los inviernos de la IA no detuvieron definitivamente el campo?
    \item \textquestiondown Qué factores permitieron que el aprendizaje profundo despegara recién después de 2012?
\end{itemize}

\begin{figure}[htpb!]
\centering
\begin{tcolorbox}[colback=white, colframe=unmsm_gold, arc=3pt, boxrule=0.8pt, top=0.3cm, bottom=0.3cm, left=0.4cm, right=0.4cm]
\centering
\small\sffamily
\textbf{Figura 1.2.} \textit{Línea de tiempo de la IA: de las primeras ideas a la era generativa}\\[0.25cm]
\resizebox{0.95\linewidth}{!}{%
\begin{tikzpicture}[>=stealth, font=\footnotesize\sffamily]
    % Eje horizontal
    \draw[->, very thick, draw=unmsm_red] (-6.8,0) -- (6.8,0);

    % Hitos
    \node[circle, fill=unmsm_red, inner sep=2.5pt] (p1) at (-5.5,0) {};
    \node[above=0.2cm of p1, align=center] {\textbf{1943--1950}\\\scriptsize McCulloch, Pitts, Turing};

    \node[circle, fill=unmsm_darkred, inner sep=2.5pt] (p2) at (-2.8,0) {};
    \node[below=0.2cm of p2, align=center] {\textbf{1956}\\\scriptsize Dartmouth};

    \node[circle, fill=unmsm_red, inner sep=2.5pt] (p3) at (-0.2,0) {};
    \node[above=0.2cm of p3, align=center] {\textbf{1974--1986}\\\scriptsize Inviernos / Backprop};

    \node[circle, fill=unmsm_darkred, inner sep=2.5pt] (p4) at (2.5,0) {};
    \node[below=0.2cm of p4, align=center] {\textbf{2012--2014}\\\scriptsize AlexNet / GAN};

    \node[circle, fill=unmsm_red, inner sep=2.5pt] (p5) at (5.5,0) {};
    \node[above=0.2cm of p5, align=center] {\textbf{2017--hoy}\\\scriptsize Transformer / GenAI};
\end{tikzpicture}%
}\\[0.2cm]
\raggedright\scriptsize \textbf{Fuente:} Elaboración propia a partir de Turing (1950), McCarthy et al. (2006) y Nasser-Eddine (2024).
\end{tcolorbox}
\end{figure}

\subsubsection{Etapa 1. Fundamentos teóricos (1943--1955)}
Antes de que existiera la palabra "inteligencia artificial" ya había ideas que la hacían pensable. En 1943, McCulloch y Pitts describieron una neurona artificial como una unidad lógica que suma señales y se activa al superar un umbral; es el germen de las redes neuronales actuales. En 1950, Turing publicó \textit{Computing machinery and intelligence}, donde propone el juego de imitación, discute las objeciones a la idea de máquinas pensantes y sugiere construir máquinas que aprendan mediante la experiencia y la educación.

En el juego de imitación, un interrogador conversa por escrito con una persona y con una máquina sin saber cuál es cuál. Si no puede distinguirlas, se considera que la máquina exhibe un comportamiento inteligente. Turing cambió así una pregunta filosófica por un criterio que se puede observar.

\begin{cajabeige}
\textbf{CONCEPTO CLAVE: Test de Turing}\\
Prueba que evalúa si una máquina puede sostener una conversación indistinguible de la de un ser humano.
\end{cajabeige}

\textbf{Personajes destacados:}
\begin{itemize}
    \item \textbf{Warren McCulloch y Walter Pitts:} Primer modelo matemático de neurona (1943).
    \item \textbf{Alan Turing:} El test de Turing y la hipótesis de máquinas que aprenden (1950).
\end{itemize}

\subsubsection{Etapa 2. Nacimiento y optimismo (1956--1973)}
En 1955, John McCarthy, Marvin Minsky, Nathaniel Rochester y Claude Shannon presentaron la propuesta de un taller de verano en Dartmouth College. Allí plantearon estudiar redes neuronales, formación de conceptos y la capacidad de las máquinas de mejorar por sí mismas. El taller de 1956 suele tomarse como el acta de nacimiento del campo. Siguieron años de entusiasmo: programas que demostraban teoremas, como el \textit{Logic Theorist} de Newell y Simon, el perceptrón de Rosenblatt (1958) y ELIZA (1966), que simulaba una conversación con reglas muy simples.

El perceptrón es un clasificador lineal: separa los datos con una recta (o un plano). En 1969, Minsky y Papert mostraron que un perceptrón de una sola capa no puede resolver problemas que no son linealmente separables, como la función XOR, lo que enfrió el interés por las redes neuronales.

Con entradas $(0,0)$, $(0,1)$, $(1,0)$ y $(1,1)$, la función XOR vale $0, 1, 1$ y $0$. Ninguna recta del plano deja los dos puntos de salida 1 de un lado y los dos de salida 0 del otro. Se necesitaban capas ocultas, y eso exigía un método para entrenarlas.

\begin{cajagris}
\textbf{CONCEPTO CLAVE: Separabilidad lineal}\\
Un conjunto de puntos es linealmente separable si una recta (o un plano) deja cada clase de un lado distinto.
\end{cajagris}

\textbf{Personajes destacados:}
\begin{itemize}
    \item \textbf{John McCarthy:} Propuso el nombre "inteligencia artificial" y el taller de Dartmouth.
    \item \textbf{Marvin Minsky, Nathaniel Rochester y Claude Shannon:} Coautores de la propuesta de 1955.
    \item \textbf{Frank Rosenblatt:} Perceptrón (1958).
    \item \textbf{Joseph Weizenbaum:} ELIZA (1966).
\end{itemize}

\subsubsection{Etapa 3. Inviernos y sistemas expertos (1974--1993)}
Las expectativas superaron lo que la tecnología podía entregar. Hacia mediados de los años setenta la financiación se redujo: es el primer invierno de la IA. En los ochenta resurgió el interés con los sistemas expertos, programas que aplicaban reglas "si\dots entonces\dots" redactadas por especialistas, y se usaron en medicina y en la industria. Pero eran caros de mantener y poco flexibles, y a fines de esa década llegó el segundo invierno. En paralelo, Rumelhart, Hinton y Williams (1986) difundieron la retropropagación, el método que hoy permite entrenar redes profundas.

Tres factores explican estos periodos:
\begin{itemize}
    \item \textbf{Promesas excesivas:} Se anunciaron logros cercanos que no llegaron.
    \item \textbf{Límites técnicos:} Poca capacidad de cómputo y escasez de datos.
    \item \textbf{Fragilidad:} Los sistemas de reglas fallaban fuera de los casos previstos y exigían un mantenimiento costoso.
\end{itemize}

\begin{figure}[htpb!]
\centering
\begin{tcolorbox}[colback=white, colframe=unmsm_gold, arc=3pt, boxrule=0.8pt, top=0.3cm, bottom=0.3cm, left=0.4cm, right=0.4cm]
\centering
\small\sffamily
\textbf{Figura 1.3.} \textit{Ciclos de expectativas e inviernos de la IA (esquema conceptual)}\\[0.25cm]
\resizebox{0.92\linewidth}{!}{%
\begin{tikzpicture}[xscale=0.95, yscale=0.8, font=\scriptsize\sffamily]
    \draw[->, thick, draw=unmsm_gray] (0,0) -- (12.5,0) node[right] {Tiempo (1950--2026)};
    \draw[->, thick, draw=unmsm_gray] (0,0) -- (0,5.5) node[above] {Expectativa e Inversión};

    \draw[thick, draw=unmsm_red, smooth] plot coordinates {
        (0.5, 1.0) (1.5, 3.8) (3.0, 1.2) (4.5, 4.2) (6.0, 1.5) (8.5, 4.5) (11.5, 5.2)
    };

    \node[anchor=south, font=\tiny\bfseries, text=unmsm_darkred] at (1.5, 3.8) {Dartmouth};
    \node[anchor=north, font=\tiny, text=unmsm_gray] at (3.0, 1.2) {Primer invierno};
    \node[anchor=south, font=\tiny\bfseries, text=unmsm_darkred] at (4.5, 4.2) {Sistemas expertos};
    \node[anchor=north, font=\tiny, text=unmsm_gray] at (6.0, 1.5) {Segundo invierno};
    \node[anchor=south, font=\tiny\bfseries, text=unmsm_darkred] at (8.5, 4.5) {AlexNet (2012)};
    \node[anchor=south, font=\tiny\bfseries, text=unmsm_red] at (11.5, 5.2) {GenAI (ChatGPT)};
\end{tikzpicture}%
}\\[0.2cm]
\raggedright\scriptsize \textbf{Fuente:} Elaboración propia a partir de Nasser-Eddine (2024). Esquema ilustrativo.
\end{tcolorbox}
\end{figure}

\begin{cajabeige}
\textbf{CONCEPTO CLAVE: Invierno de la IA}\\
Periodo en que bajan la financiación y el interés en la IA porque los resultados no cumplen las expectativas generadas.
\end{cajabeige}

\subsubsection{Etapa 4. Aprendizaje profundo (1997--2016)}
Tres factores cambiaron el panorama, resumidos en la Tabla 1.5. En 1997, Deep Blue derrotó a Kaspárov en ajedrez, aunque lo hizo sobre todo con búsqueda masiva y no con aprendizaje.

\begin{table}[htpb!]
\centering
\small\sffamily
\caption{\textbf{Tabla 1.5.} \textit{Por qué despegó el aprendizaje profundo}}
\vspace{0.15cm}
\begin{tabularx}{\textwidth}{>{\bfseries}p{2.5cm} >{\raggedright\arraybackslash}p{6.5cm} >{\raggedright\arraybackslash}X}
\toprule
Factor & Qué cambió & Efecto \\
\midrule
Datos & Internet y los dispositivos generaron enormes colecciones de texto, imágenes y audio. & Los modelos pudieron aprender patrones más ricos. \\
\midrule
Cómputo & Los procesadores gráficos (GPU) permitieron entrenar redes con millones de parámetros. & El entrenamiento pasó de meses a días. \\
\midrule
Algoritmos & Mejores funciones de activación, inicialización y técnicas de entrenamiento. & Las redes de muchas capas dejaron de ser inestables. \\
\bottomrule
\end{tabularx}
\vspace{0.15cm}
\raggedright\scriptsize \textbf{Fuente:} Elaboración propia a partir de LeCun, Bengio y Hinton (2015).
\end{table}

El punto de quiebre llegó en 2012, cuando una red convolucional profunda, AlexNet, ganó el concurso ImageNet con una ventaja que sorprendió a la comunidad. En 2014, Goodfellow et al. introdujeron las redes generativas antagónicas; en 2016, el programa AlphaGo venció a un campeón de Go; y en 2015 LeCun, Bengio y Hinton resumieron en \textit{Nature} cómo el aprendizaje profundo aprende representaciones a partir de datos masivos.

\textbf{Personajes destacados:}
\begin{itemize}
    \item \textbf{Yann LeCun, Yoshua Bengio y Geoffrey Hinton:} Impulsores del aprendizaje profundo.
    \item \textbf{Ian Goodfellow:} Redes generativas antagónicas, GAN (2014).
\end{itemize}

\subsubsection{Etapa 5. Era generativa (2017--hoy)}
En 2017, un equipo de Google presentó el Transformer, una arquitectura basada en "atención" que procesa secuencias en paralelo sin recurrir a redes recurrentes ni convolucionales, con mayor precisión y menor tiempo de entrenamiento en traducción automática. Sobre esa base creció la familia GPT. GPT-3, con 175 mil millones de parámetros, mostró que un modelo suficientemente grande puede resolver tareas nuevas con solo unos pocos ejemplos en la instrucción, sin reentrenarse. Ese mismo año, Ho, Jain y Abbeel (2020) propusieron los modelos de difusión, que generan imágenes de alta calidad eliminando ruido paso a paso.

En 2021 se acuñó el término modelos fundacionales para describir sistemas entrenados con enormes cantidades de datos y adaptables a muchas tareas. En 2022, Ouyang et al. mostraron que la retroalimentación humana vuelve a los modelos más útiles y menos tóxicos que GPT-3, y poco después los asistentes conversacionales llegaron al público masivo. Zhao et al. (2023) ofrecen un panorama de cómo se entrenan y evalúan estos grandes modelos de lenguaje.

La literatura reciente muestra además que el campo todavía busca consensos. García-Peñalvo y Vázquez-Ingelmo (2023) observan una diferencia entre la percepción del público, centrada en la creación de contenido, y la de la comunidad científica, que presta más atención a los modelos y técnicas. Nasser-Eddine (2024) añade que la IA generativa plantea desafíos de privacidad, sesgos, seguridad y propiedad intelectual.

\textbf{Personajes destacados:}
\begin{itemize}
    \item \textbf{Ashish Vaswani y colegas:} El Transformer (2017).
    \item \textbf{Tom Brown y colegas (OpenAI):} GPT-3 y el aprendizaje con pocos ejemplos (2020).
    \item \textbf{Jonathan Ho, Ajay Jain y Pieter Abbeel:} Modelos de difusión probabilística (2020).
    \item \textbf{Long Ouyang y colegas:} Entrenamiento con retroalimentación humana, RLHF (2022).
\end{itemize}

\vspace{0.3cm}

\subsection{Cómo aprende una máquina}

En la programación tradicional, una persona escribe las reglas y la computadora las aplica a los datos para dar una respuesta. En el aprendizaje automático (\textit{machine learning}) se invierte la lógica: se entregan datos y ejemplos de respuestas, y el sistema descubre las reglas por sí mismo. Esa diferencia explica por qué estos sistemas resuelven problemas para los que nadie sabría escribir reglas a mano, como reconocer un rostro o seguir una conversación.

\subsubsection{Aprendizaje automático y sus tipos}
El aprendizaje automático permite que un sistema mejore su desempeño en una tarea a partir de la experiencia, es decir, de los datos. Según el tipo de datos disponibles y la forma de "corregir" al sistema, se distinguen varios tipos de aprendizaje (Tabla 1.6).

\begin{cajagris}
\textbf{CONCEPTO CLAVE: Aprendizaje Automático}\\
Conjunto de métodos que permiten a un sistema aprender patrones a partir de datos para hacer predicciones o tomar decisiones.
\end{cajagris}

\begin{table}[htpb!]
\centering
\small\sffamily
\caption{\textbf{Tabla 1.6.} \textit{Principales tipos de aprendizaje automático}}
\vspace{0.15cm}
\begin{tabularx}{\textwidth}{>{\bfseries}p{3.0cm} >{\raggedright\arraybackslash}p{3.8cm} >{\raggedright\arraybackslash}p{4.2cm} >{\raggedright\arraybackslash}X}
\toprule
Tipo & Datos que usa & Cómo aprende & Ejemplo \\
\midrule
Supervisado & Ejemplos con respuesta conocida (etiquetas). & Reduce la diferencia entre su predicción y la etiqueta. & Clasificar correos como spam o no spam. \\
\midrule
No supervisado & Datos sin etiquetas. & Busca estructura: grupos, similitudes, patrones. & Agrupar clientes con hábitos parecidos. \\
\midrule
Por refuerzo & Interacción con un entorno. & Prueba acciones y recibe recompensas o castigos. & Un programa que aprende a jugar ajedrez o Go. \\
\midrule
Autosupervisado & Datos sin etiquetar de forma manual. & Genera sus propias etiquetas a partir de los datos. & Predecir la palabra siguiente en un texto. \\
\bottomrule
\end{tabularx}
\vspace{0.15cm}
\raggedright\scriptsize \textbf{Fuente:} Elaboración propia.
\end{table}

El aprendizaje autosupervisado es la base de los modelos de lenguaje: el propio texto proporciona la "respuesta correcta", lo que permite entrenar con cantidades enormes de datos sin etiquetarlos a mano.

\subsubsection{Redes neuronales y aprendizaje profundo}
Una red neuronal artificial está formada por capas de unidades sencillas, las "neuronas". Cada neurona recibe números, los multiplica por pesos, los suma, añade un sesgo y aplica una función de activación:
\begin{equation}
y = f\left( w_1 x_1 + w_2 x_2 + \dots + w_n x_n + b \right)
\end{equation}
La función de activación $f$ introduce no linealidad; sin ella, apilar capas equivaldría a una sola operación lineal. La Tabla 1.7 presenta las más usadas.

\begin{table}[htpb!]
\centering
\small\sffamily
\caption{\textbf{Tabla 1.7.} \textit{Funciones de activación frecuentes en redes neuronales}}
\vspace{0.15cm}
\begin{tabularx}{\textwidth}{>{\bfseries}p{3.2cm} >{\raggedright\arraybackslash}p{4.0cm} >{\raggedright\arraybackslash}p{3.0cm} >{\raggedright\arraybackslash}X}
\toprule
Función & Fórmula & Rango & Ejemplo \\
\midrule
ReLU & $f(z) = \max(0, z)$ & $[0, \infty)$ & $f(-2) = 0$; $f(3) = 3$ \\
\midrule
Sigmoide & $\sigma(z) = \frac{1}{1 + e^{-z}}$ & $(0, 1)$ & $\sigma(0) = 0.5$ \\
\midrule
Tangente hiperbólica & $\tanh(z) = \frac{e^z - e^{-z}}{e^z + e^{-z}}$ & $(-1, 1)$ & $\tanh(0) = 0$ \\
\bottomrule
\end{tabularx}
\vspace{0.15cm}
\raggedright\scriptsize \textbf{Fuente:} Elaboración propia.
\end{table}

Al apilar muchas capas se obtiene una red profunda, de ahí el nombre aprendizaje profundo. LeCun, Bengio y Hinton (2015) explican que cada capa transforma la representación de la anterior en otra algo más abstracta: en una imagen, los píxeles dan lugar a bordes, luego a formas y finalmente a objetos.

\begin{figure}[htpb!]
\centering
\begin{tcolorbox}[colback=white, colframe=unmsm_gold, arc=3pt, boxrule=0.8pt, top=0.3cm, bottom=0.3cm, left=0.4cm, right=0.4cm]
\centering
\small\sffamily
\textbf{Figura 1.4.} \textit{Estructura de una red neuronal con dos capas ocultas}\\[0.25cm]
\resizebox{0.92\linewidth}{!}{%
\begin{tikzpicture}[>=stealth, font=\footnotesize\sffamily]
    % Capa Entrada
    \foreach \y in {1,2,3} {
        \node[circle, draw=unmsm_red, fill=bg_marfil, line width=1pt] (in\y) at (0, 2-\y*1.2) {$x_\y$};
    }
    \node[above=0.2cm of in1] {\textbf{Entrada}};

    % Capa Oculta 1
    \foreach \y in {1,2,3,4} {
        \node[circle, draw=unmsm_darkred, fill=bg_arena, line width=1pt] (h1\y) at (3.5, 2.5-\y*1.1) {$h^{(1)}_\y$};
    }
    \node[above=0.2cm of h11] {\textbf{Capa Oculta 1}};

    % Capa Oculta 2
    \foreach \y in {1,2,3,4} {
        \node[circle, draw=unmsm_darkred, fill=bg_arena, line width=1pt] (h2\y) at (7.0, 2.5-\y*1.1) {$h^{(2)}_\y$};
    }
    \node[above=0.2cm of h21] {\textbf{Capa Oculta 2}};

    % Capa Salida
    \node[circle, draw=unmsm_red, fill=bg_marfil, line width=1pt] (out) at (10.5, 0) {$\hat{y}$};
    \node[above=0.2cm of out] {\textbf{Salida}};

    % Conexiones
    \foreach \i in {1,2,3} {
        \foreach \j in {1,2,3,4} {
            \draw[->, draw=unmsm_gray!50] (in\i) -- (h1\j);
        }
    }
    \foreach \i in {1,2,3,4} {
        \foreach \j in {1,2,3,4} {
            \draw[->, draw=unmsm_gray!50] (h1\i) -- (h2\j);
        }
    }
    \foreach \i in {1,2,3,4} {
        \draw[->, draw=unmsm_darkred] (h2\i) -- (out);
    }
\end{tikzpicture}%
}\\[0.2cm]
\raggedright\scriptsize \textbf{Fuente:} Elaboración propia a partir de LeCun, Bengio y Hinton (2015).
\end{tcolorbox}
\end{figure}

\subsubsection{Dos arquitecturas clásicas}
\begin{itemize}
    \item \textbf{Redes convolucionales (CNN):} Aplican filtros que recorren la imagen y detectan bordes, texturas y formas; son útiles para imágenes y sonido.
    \item \textbf{Redes recurrentes (RNN):} Procesan los datos paso a paso conservando una memoria de lo anterior; se usaron para texto y voz, pero les cuesta manejar secuencias largas, una limitación que el Transformer buscó superar.
\end{itemize}

\subsubsection{Retropropagación y descenso del gradiente}
\textquestiondown Cómo se ajustan los pesos? Con el algoritmo de retropropagación: la red hace una predicción, se mide cuánto se equivoca mediante una función de pérdida y el error se propaga hacia atrás para calcular cuánto contribuyó cada peso. Luego cada peso se corrige un poco en la dirección que reduce el error, un método llamado descenso del gradiente:
\begin{equation}
w \leftarrow w - \eta \cdot \frac{\partial L}{\partial w}
\end{equation}
donde $L$ es la pérdida y $\eta$ (\textit{eta}) es la tasa de aprendizaje, un número pequeño que controla el tamaño de cada paso.

\subsubsection{Datos, entrenamiento y modelos}
Tres piezas aparecen siempre que hablamos de aprendizaje automático:
\begin{itemize}
    \item \textbf{Datos:} Son la materia prima: textos, imágenes, tablas o audios. Su cantidad y su calidad condicionan lo que el sistema puede aprender; si los datos son sesgados o incompletos, el modelo heredará esas fallas.
    \item \textbf{Entrenamiento:} Es el proceso de ajustar los parámetros para minimizar la pérdida sobre los datos. Se repite muchas veces sobre los mismos ejemplos; cada recorrido completo se llama época.
    \item \textbf{Modelo:} Es el resultado: una función con parámetros ya ajustados. Los grandes modelos de lenguaje tienen miles de millones de ellos (GPT-3 tiene 175 mil millones, según Brown et al., 2020). Al usarlo con datos nuevos se habla de inferencia.
\end{itemize}

\begin{figure}[htpb!]
\centering
\begin{tcolorbox}[colback=white, colframe=unmsm_gold, arc=3pt, boxrule=0.8pt, top=0.3cm, bottom=0.3cm, left=0.4cm, right=0.4cm]
\centering
\small\sffamily
\textbf{Figura 1.5.} \textit{Del conjunto de datos al modelo en uso}\\[0.25cm]
\resizebox{0.95\linewidth}{!}{%
\begin{tikzpicture}[>=stealth, font=\footnotesize\sffamily]
    \node [rectangle, draw=unmsm_red, fill=bg_marfil, rounded corners=2pt, text width=2.4cm, text centered] (d) at (-4.5, 0) {\textbf{1. Datos}\\Texto, imágenes, tablas};
    \node [rectangle, draw=unmsm_darkred, fill=bg_arena, rounded corners=2pt, text width=2.6cm, text centered] (t) at (-1.5, 0) {\textbf{2. Entrenamiento}\\Minimizar error de pérdida};
    \node [rectangle, draw=unmsm_red, fill=bg_marfil, rounded corners=2pt, text width=2.4cm, text centered] (m) at (1.5, 0) {\textbf{3. Modelo}\\Pesos ajustados};
    \node [rectangle, draw=unmsm_darkred, fill=white, rounded corners=2pt, text width=2.4cm, text centered] (i) at (4.5, 0) {\textbf{4. Inferencia}\\Datos nuevos};

    \draw [->, thick, draw=unmsm_red] (d) -- (t);
    \draw [->, thick, draw=unmsm_darkred] (t) -- (m);
    \draw [->, thick, draw=unmsm_red] (m) -- (i);
\end{tikzpicture}%
}\\[0.2cm]
\raggedright\scriptsize \textbf{Fuente:} Elaboración propia.
\end{tcolorbox}
\end{figure}

Para saber si un modelo realmente aprendió, y no solo memorizó, los datos se dividen habitualmente en tres grupos: entrenamiento (70\%), validación (15\%) y prueba (15\%). Si rinde muy bien con los datos de entrenamiento pero mal con los de prueba, hay sobreajuste.

\begin{cajagris}
\textbf{CONCEPTO CLAVE: Sobreajuste (Overfitting)}\\
Ocurre cuando el modelo memoriza los ejemplos de entrenamiento y pierde capacidad de generalizar a datos nuevos.
\end{cajagris}

\begin{figure}[htpb!]
\centering
\begin{tcolorbox}[colback=white, colframe=unmsm_gold, arc=3pt, boxrule=0.8pt, top=0.3cm, bottom=0.3cm, left=0.4cm, right=0.4cm]
\centering
\small\sffamily
\textbf{Figura 1.6.} \textit{Error de entrenamiento y de validación según la complejidad del modelo}\\[0.25cm]
\resizebox{0.90\linewidth}{!}{%
\begin{tikzpicture}[font=\scriptsize\sffamily]
    \draw[->, thick, draw=unmsm_gray] (0,0) -- (7.5,0) node[right] {Complejidad / Épocas $\to$};
    \draw[->, thick, draw=unmsm_gray] (0,0) -- (0,4.5) node[above] {Error};

    % Curva train
    \draw[thick, draw=unmsm_red] (0.5, 4.0) to[out=-60, in=175] (7.0, 0.4);
    \node[text=unmsm_red, font=\tiny\bfseries] at (6.0, 0.8) {Error entrenamiento};

    % Curva val
    \draw[thick, draw=unmsm_darkred] (0.5, 4.2) to[out=-55, in=180] (3.5, 1.6) to[out=0, in=-120] (7.0, 3.8);
    \node[text=unmsm_darkred, font=\tiny\bfseries] at (6.0, 3.4) {Error validación};

    % Punto optimo
    \draw[dashed, draw=gray] (3.5, 0) -- (3.5, 1.6);
    \node[below, font=\tiny] at (3.5, 0) {Punto óptimo};
\end{tikzpicture}%
}\\[0.2cm]
\raggedright\scriptsize \textbf{Fuente:} Elaboración propia. Esquema ilustrativo.
\end{tcolorbox}
\end{figure}

Hay que distinguir dos tipos de números: los \textbf{parámetros} (pesos y sesgos) los aprende el modelo durante el entrenamiento. Los \textbf{hiperparámetros} (como la tasa de aprendizaje $\eta$, el número de capas o de épocas) los decide la persona que diseña y entrena el modelo.

\subsubsection{Ejemplo: un paso de entrenamiento}
Consideremos el modelo más simple, $\hat{y} = w \cdot x$, con pérdida cuadrática $L = (\hat{y} - y)^2$. Para el dato $x = 2$ la respuesta correcta es $y = 6$, el peso inicial es $w = 1$ y la tasa de aprendizaje es $\eta = 0.05$. La Tabla 1.8 muestra el cálculo paso a paso.

\begin{table}[htpb!]
\centering
\small\sffamily
\caption{\textbf{Tabla 1.8.} \textit{Un paso de descenso del gradiente}}
\vspace{0.15cm}
\begin{tabularx}{\textwidth}{>{\bfseries}p{3.2cm} >{\raggedright\arraybackslash}p{6.0cm} >{\raggedright\arraybackslash}X}
\toprule
Paso & Cálculo & Resultado \\
\midrule
1. Predicción & $\hat{y} = 1 \cdot 2$ & $\hat{y} = 2$ \\
\midrule
2. Pérdida & $L = (2 - 6)^2$ & $L = 16$ \\
\midrule
3. Gradiente & $\frac{\partial L}{\partial w} = 2(\hat{y} - y) \cdot x = 2(-4)(2)$ & $-16$ \\
\midrule
4. Actualización & $w = 1 - 0.05 \cdot (-16)$ & $w = 1.8$ \\
\midrule
5. Nueva pérdida & $\hat{y} = 1.8 \cdot 2 = 3.6$; $L = (3.6 - 6)^2$ & $L = 5.76$ \\
\bottomrule
\end{tabularx}
\vspace{0.15cm}
\raggedright\scriptsize \textbf{Fuente:} Elaboración propia.
\end{table}

Con un solo paso, la pérdida bajó de 16 a 5.76. Repitiendo el proceso, $w$ se acerca a 3, el valor que hace que la predicción coincida exactamente con la respuesta. Un modelo real hace exactamente esto, pero con miles de millones de parámetros y millones de ejemplos.

\subsubsection{Cómo medir si el modelo funciona}
Supongamos un filtro de spam evaluado con 100 correos de prueba, 40 de ellos spam. La Tabla 1.9 resume sus aciertos y errores (datos hipotéticos).

\begin{table}[htpb!]
\centering
\small\sffamily
\caption{\textbf{Tabla 1.9.} \textit{Resultados de un filtro de spam sobre datos de prueba}}
\vspace{0.15cm}
\begin{tabularx}{\textwidth}{>{\bfseries}p{4.0cm} >{\raggedright\arraybackslash}p{4.0cm} >{\raggedright\arraybackslash}X}
\toprule
Realidad $\backslash$ Predicción & Predice "spam" & Predice "no spam" \\
\midrule
Es spam (40) & 34 aciertos & 6 errores \\
\midrule
No es spam (60) & 3 errores & 57 aciertos \\
\bottomrule
\end{tabularx}
\vspace{0.15cm}
\raggedright\scriptsize \textbf{Fuente:} Elaboración propia.
\end{table}

La exactitud (\textit{accuracy}) es la proporción de aciertos: $(34 + 57) / 100 = 91\%$. Pero también interesa cuántos spam detecta de los que existen (sensibilidad o \textit{recall}): $34 / 40 = 85\%$. Según el problema, importa más un tipo de error u otro; por eso no basta con un solo número.

\subsubsection{De las redes a los modelos de lenguaje}
Todo lo anterior se aplica a los modelos de lenguaje. Un modelo como GPT-3 es una red neuronal muy grande, entrenada de forma autosupervisada sobre grandes cantidades de texto para predecir la palabra siguiente. Su "conocimiento" está codificado en miles de millones de pesos ajustados con el mismo principio de pérdida y descenso del gradiente. En la sección 1.4 veremos qué cambia cuando el objetivo ya no es clasificar, sino generar.

\vspace{0.3cm}

\subsection{La IA generativa}

La IA generativa produce contenido nuevo, como texto, imágenes, audio o código, a partir de los patrones aprendidos de los datos. No es una rama aparte, sino un resultado del aprendizaje profundo y de la escala de datos y cómputo vistos en la sección anterior.

\subsubsection{Clasificar o predecir frente a generar}
Un modelo discriminativo aprende a decidir entre opciones: dada una entrada, responde "\textquestiondown qué es?" o "\textquestiondown cuánto valdrá?". Un modelo generativo aprende cómo son los datos mismos y puede producir ejemplos nuevos que se parecen a los originales (Tabla 1.10 y Tabla 1.11).

\begin{table}[htpb!]
\centering
\small\sffamily
\caption{\textbf{Tabla 1.10.} \textit{Diferencias entre clasificar o predecir y generar}}
\vspace{0.15cm}
\begin{tabularx}{\textwidth}{>{\bfseries}p{3.2cm} >{\raggedright\arraybackslash}p{5.5cm} >{\raggedright\arraybackslash}X}
\toprule
Criterio & IA que clasifica o predice & IA que genera \\
\midrule
Pregunta que responde & "\textquestiondown Esta imagen contiene un gato?" & "Dibuja un gato dormido sobre un libro." \\
\midrule
Resultado & Una etiqueta o un número. & Contenido nuevo: texto, imagen, audio, código. \\
\midrule
Ejemplo típico & Filtro de spam, diagnóstico de imágenes. & Asistente conversacional, generador de imágenes. \\
\midrule
Cómo se evalúa & Aciertos frente a una respuesta conocida. & Calidad, coherencia y veracidad; a menudo hay varias respuestas válidas. \\
\bottomrule
\end{tabularx}
\vspace{0.15cm}
\raggedright\scriptsize \textbf{Fuente:} Elaboración propia a partir de Feuerriegel et al. (2024).
\end{table}

\begin{table}[htpb!]
\centering
\small\sffamily
\caption{\textbf{Tabla 1.11.} \textit{Tipos de contenido que genera la IA}}
\vspace{0.15cm}
\begin{tabularx}{\textwidth}{>{\bfseries}p{2.5cm} >{\raggedright\arraybackslash}p{5.5cm} >{\raggedright\arraybackslash}X}
\toprule
Contenido & Enfoque habitual & Usos frecuentes \\
\midrule
Texto & Modelos de lenguaje basados en Transformer. & Redactar, resumir, traducir, responder preguntas. \\
\midrule
Imagen & Modelos de difusión y redes generativas antagónicas. & Ilustraciones, edición de imágenes, datos sintéticos. \\
\midrule
Audio & Modelos generativos de voz y sonido. & Síntesis de voz, música, restauración de audio. \\
\midrule
Código & Modelos de lenguaje especializados en programación. & Completar funciones, explicar errores, generar pruebas. \\
\bottomrule
\end{tabularx}
\vspace{0.15cm}
\raggedright\scriptsize \textbf{Fuente:} Elaboración propia a partir de Feuerriegel et al. (2024) y García-Peñalvo y Vázquez-Ingelmo (2023).
\end{table}

\subsubsection{Idea básica de los modelos de lenguaje}
Un modelo de lenguaje divide el texto en pequeñas piezas llamadas tokens (palabras o fragmentos de palabra) y aprende a estimar la probabilidad del siguiente token dado todo lo anterior. Para escribir una respuesta, elige un token, lo añade al texto y repite el proceso (Figura 1.7).

\begin{figure}[htpb!]
\centering
\begin{tcolorbox}[colback=white, colframe=unmsm_gold, arc=3pt, boxrule=0.8pt, top=0.3cm, bottom=0.3cm, left=0.4cm, right=0.4cm]
\centering
\small\sffamily
\textbf{Figura 1.7.} \textit{Predicción del siguiente token (valores ilustrativos)}\\[0.25cm]
\resizebox{0.90\linewidth}{!}{%
\begin{tikzpicture}[font=\small\sffamily]
    \node[anchor=west] at (0, 2.5) {\textbf{Entrada (prompt):} "La capital del Perú es \dots" $\to$ \textquestiondown qué token sigue?};

    \draw[fill=unmsm_red, draw=unmsm_darkred] (0, 1.8) rectangle (6.88, 2.2);
    \node[anchor=west] at (7.1, 2.0) {\textbf{Lima} (0.86)};

    \draw[fill=bg_arena, draw=unmsm_gray] (0, 1.2) rectangle (0.4, 1.6);
    \node[anchor=west] at (0.6, 1.4) {Luna (0.05)};

    \draw[fill=bg_arena, draw=unmsm_gray] (0, 0.6) rectangle (0.32, 1.0);
    \node[anchor=west] at (0.5, 0.8) {Arequipa (0.04)};

    \draw[fill=bg_arena, draw=unmsm_gray] (0, 0.0) rectangle (0.24, 0.4);
    \node[anchor=west] at (0.4, 0.2) {Cusco (0.03)};
\end{tikzpicture}%
}\\[0.2cm]
\raggedright\scriptsize \textbf{Fuente:} Elaboración propia.
\end{tcolorbox}
\end{figure}

Tres ingredientes explican su capacidad actual:
\begin{enumerate}
    \item \textbf{Arquitectura Transformer:} Su mecanismo de autoatención permite relacionar cada palabra con las demás del texto y procesarlas en paralelo.
    \item \textbf{Escala:} Con más datos y parámetros, los modelos pueden aprender nuevas tareas a partir de pocos ejemplos escritos en la instrucción.
    \item \textbf{Alineación con las personas:} La retroalimentación humana (RLHF) enseña al modelo a seguir instrucciones de forma más útil y segura.
\end{enumerate}

Bommasani et al. (2021) llaman \textbf{modelos fundacionales} a este tipo de sistemas y subrayan dos rasgos: la aparición de capacidades nuevas con la escala y la reutilización de un mismo modelo para muchas aplicaciones.

\subsubsection{Cómo se alinea un modelo con las instrucciones humanas}
Un modelo entrenado solo para predecir texto no sigue necesariamente instrucciones. Ouyang et al. (2022) mostraron que añadir retroalimentación humana mejora el resultado, siguiendo tres pasos (Figura 1.8). Con este método, el modelo InstructGPT produjo respuestas más útiles, precisas y menos tóxicas que GPT-3, y se comprobó que el tamaño por sí solo no garantiza un mejor comportamiento.

\begin{figure}[htpb!]
\centering
\begin{tcolorbox}[colback=white, colframe=unmsm_gold, arc=3pt, boxrule=0.8pt, top=0.3cm, bottom=0.3cm, left=0.4cm, right=0.4cm]
\centering
\small\sffamily
\textbf{Figura 1.8.} \textit{Aprendizaje por refuerzo a partir de retroalimentación humana (RLHF)}\\[0.25cm]
\resizebox{0.95\linewidth}{!}{%
\begin{tikzpicture}[>=stealth, font=\footnotesize\sffamily]
    \node [rectangle, draw=unmsm_red, fill=bg_marfil, rounded corners=2pt, text width=3.4cm, text centered] (f1) at (-4.2, 0) {\textbf{1. Demostraciones}\\Personas escriben respuestas modelo a distintas instrucciones.};
    \node [rectangle, draw=unmsm_darkred, fill=bg_arena, rounded corners=2pt, text width=3.4cm, text centered] (f2) at (0, 0) {\textbf{2. Comparaciones}\\Personas ordenan varias respuestas. Se entrena modelo de recompensa.};
    \node [rectangle, draw=unmsm_red, fill=bg_marfil, rounded corners=2pt, text width=3.4cm, text centered] (f3) at (4.2, 0) {\textbf{3. Refuerzo}\\El modelo se optimiza para obtener mayor recompensa.};

    \draw [->, thick, draw=unmsm_red] (f1) -- (f2);
    \draw [->, thick, draw=unmsm_darkred] (f2) -- (f3);
\end{tikzpicture}%
}\\[0.2cm]
\raggedright\scriptsize \textbf{Fuente:} Elaboración propia a partir de Ouyang et al. (2022).
\end{tcolorbox}
\end{figure}

\subsubsection{Cómo pedir bien: las instrucciones o prompts}
Como estos modelos responden a lo que se les escribe, la calidad de la instrucción importa. Brown et al. (2020) mostraron que basta con incluir unos pocos ejemplos dentro de la propia instrucción (\textit{few-shot}), como en este caso:
\begin{quote}
\textit{Traduce al inglés. Ejemplo 1: "gato" $\to$ cat. Ejemplo 2: "casa" $\to$ house. Ahora tú: "libro" $\to$}
\end{quote}
Tres recomendaciones prácticas:
\begin{itemize}
    \item \textbf{Sé específico:} Indica la tarea, el contexto y el formato de respuesta que esperas.
    \item \textbf{Da ejemplos:} Especialmente cuando el formato o estilo de salida importa.
    \item \textbf{Pide los pasos o las fuentes y verifícalos:} Una respuesta convincente no es necesariamente correcta.
\end{itemize}

\subsubsection{Otras familias generativas: GAN y difusión}
Los modelos de lenguaje no son los únicos. Para imágenes y otros datos se usan principalmente dos enfoques: las redes generativas antagónicas (donde un generador y un discriminador compiten) y los modelos de difusión (que aprenden a revertir un proceso gradual de adición de ruido). La Tabla 1.12 compara estas tres familias.

\begin{table}[htpb!]
\centering
\small\sffamily
\caption{\textbf{Tabla 1.12.} \textit{Comparación de arquitecturas generativas}}
\vspace{0.15cm}
\begin{tabularx}{\textwidth}{>{\bfseries}p{2.6cm} >{\raggedright\arraybackslash}p{4.2cm} >{\raggedright\arraybackslash}p{4.0cm} >{\raggedright\arraybackslash}X}
\toprule
Arquitectura & Idea central & Ventaja & Limitación \\
\midrule
Transformer & Atención entre todos los elementos de la secuencia. & Paralelismo y buen desempeño en lenguaje. & Puede alucinar; consume muchos recursos computacionales. \\
\midrule
GAN & Generador contra discriminador. & Imágenes realistas y rápidas de generar. & Entrenamiento inestable; poca variedad (colapso de modos). \\
\midrule
Difusión & Eliminar ruido de forma iterativa. & Alta calidad, comparable a las GAN en algunas pruebas. & Generación en varios pasos (más lenta). \\
\bottomrule
\end{tabularx}
\vspace{0.15cm}
\raggedright\scriptsize \textbf{Fuente:} Elaboración propia a partir de Goodfellow et al. (2014), de la Torre (2023), Ho et al. (2020) y Vaswani et al. (2017).
\end{table}

\subsubsection{Capacidades, límites y riesgos}
La potencia de estos modelos viene acompañada de problemas que la literatura revisada documenta con insistencia (Tabla 1.13).

\begin{table}[htpb!]
\centering
\small\sffamily
\caption{\textbf{Tabla 1.13.} \textit{Principales desafíos y riesgos de la IA generativa}}
\vspace{0.15cm}
\begin{tabularx}{\textwidth}{>{\bfseries}p{3.5cm} >{\raggedright\arraybackslash}p{7.5cm} >{\raggedright\arraybackslash}X}
\toprule
Desafío & En qué consiste & Fuente \\
\midrule
Alucinaciones & Contenido que contradice la fuente (internas) o que no puede verificarse con ella (externas). & Ji et al. (2023) \\
\midrule
Falta de comprensión & Generar texto coherente no implica comprender su significado real. & Bender et al. (2021) \\
\midrule
Sesgos y desinformación & Reproducción de prejuicios presentes en los datos de entrenamiento. & Bommasani et al. (2021) \\
\midrule
Costos ambientales & Entrenar modelos cada vez mayores exige mucha energía y recursos de refrigeración. & Bender et al. (2021) \\
\midrule
Dependencia acrítica & Aceptar las respuestas sin verificarlas ni entender cómo se producen. & Artopoulos y Lliteras (2024) \\
\bottomrule
\end{tabularx}
\vspace{0.15cm}
\raggedright\scriptsize \textbf{Fuente:} Elaboración propia a partir de la literatura citada.
\end{table}

Estos límites no invalidan la tecnología, pero exigen supervisión humana, colaboración entre personas e IA y un uso responsable. Hay un contraste de fondo: mientras el rendimiento técnico crece con la escala, los problemas de base, como la comprensión semántica y las alucinaciones, persisten.

\vspace{0.3cm}

\subsection{Caso o aplicación: un asistente generativo en un curso universitario}

El siguiente caso es ilustrativo: describe una situación plausible construida a partir de los problemas que documenta la literatura, no un hecho reportado en una de las fuentes.

\subsubsection*{Contexto:}
Una estudiante de primer ciclo prepara un trabajo para un curso de matemática. Para avanzar más rápido, le pide a un asistente conversacional que resuelva un ejercicio y le indique "dos referencias bibliográficas que respalden el método". Sánchez y Carbajal (2023) y García-Peñalvo (2024) señalan que estas herramientas ofrecen oportunidades reales en la educación universitaria, pero también riesgos si se usan sin criterio.

\subsubsection*{Análisis:}
El asistente entrega una solución bien redactada y dos referencias con autores, año y revista. Al revisarlas, la estudiante descubre que una de las referencias no existe y que un paso del procedimiento contiene un error de cálculo. Es el comportamiento que describen Ji et al. (2023): una alucinación externa (la cita no puede verificarse) y otra interna (el resultado contradice el propio planteamiento). García-Peñalvo, Llorens-Largo y Vidal (2024) explican la causa de fondo: el modelo predice texto probable, no "consulta" una fuente verificada. Como recuerdan Bender et al. (2021), la fluidez del texto no garantiza que el sistema comprenda su contenido.

La Tabla 1.14 resume las señales de alerta que la estudiante pudo haber detectado antes de confiar en la salida.

\begin{table}[htpb!]
\centering
\small\sffamily
\caption{\textbf{Tabla 1.14.} \textit{Señales de alerta ante respuestas generadas por IA}}
\vspace{0.15cm}
\begin{tabularx}{\textwidth}{>{\bfseries}p{5.5cm} >{\raggedright\arraybackslash}X}
\toprule
Señal de alerta & Qué hacer \\
\midrule
Una cita sin enlace o DOI comprobable. & Buscar el título en una base de datos académica antes de usarla. \\
\midrule
Cifras o afirmaciones sin fuente. & Contrastarlas con un texto de referencia. \\
\midrule
Un resultado matemático que se presenta sin justificación. & Rehacer el cálculo por otro método o pedir los pasos desglosados. \\
\midrule
La respuesta cambia al repetir la pregunta. & Tratarla con cautela: el modelo genera probabilísticamente, no consulta una base fija. \\
\midrule
Seguridad excesiva ante un tema ambiguo. & Pedir que justifique y buscar una segunda fuente independiente. \\
\bottomrule
\end{tabularx}
\vspace{0.15cm}
\raggedright\scriptsize \textbf{Fuente:} Elaboración propia a partir de Ji et al. (2023) y Artopoulos y Lliteras (2024).
\end{table}

\subsubsection*{Interpretación y Resultados:}
La herramienta fue útil como punto de partida, pero no como fuente de verdad. La estudiante tuvo que verificar cada afirmación, comprobar las referencias en bases de datos académicas y rehacer el cálculo. Esto refleja lo que proponen Artopoulos y Lliteras (2024): desarrollar alfabetización crítica en IA, entender cómo funcionan estos sistemas y evaluar sus resultados. También aplica la recomendación de García-Peñalvo (2024) de declarar el uso de IA: qué herramienta se usó, qué se pidió y cómo se verificó el resultado (Figura 1.9).

\begin{figure}[htpb!]
\centering
\begin{tcolorbox}[colback=white, colframe=unmsm_gold, arc=3pt, boxrule=0.8pt, top=0.3cm, bottom=0.3cm, left=0.4cm, right=0.4cm]
\centering
\small\sffamily
\textbf{Figura 1.9.} \textit{Flujo de uso responsable de un asistente generativo}\\[0.25cm]
\resizebox{0.95\linewidth}{!}{%
\begin{tikzpicture}[>=stealth, font=\footnotesize\sffamily]
    \node [rectangle, draw=unmsm_red, fill=bg_marfil, rounded corners=2pt, text width=2.4cm, text centered] (p1) at (-4.5, 0) {\textbf{1. Pedir}\\Tarea con contexto y fuentes.};
    \node [rectangle, draw=unmsm_darkred, fill=bg_arena, rounded corners=2pt, text width=2.4cm, text centered] (p2) at (-1.5, 0) {\textbf{2. Revisar}\\Releer y rehacer cálculos.};
    \node [rectangle, draw=unmsm_red, fill=bg_marfil, rounded corners=2pt, text width=2.4cm, text centered] (p3) at (1.5, 0) {\textbf{3. Contrastar}\\Comprobar citas en bases reales.};
    \node [rectangle, draw=unmsm_darkred, fill=white, rounded corners=2pt, text width=2.4cm, text centered] (p4) at (4.5, 0) {\textbf{4. Declarar}\\Indicar herramienta y verificación.};

    \draw [->, thick, draw=unmsm_red] (p1) -- (p2);
    \draw [->, thick, draw=unmsm_darkred] (p2) -- (p3);
    \draw [->, thick, draw=unmsm_red] (p3) -- (p4);
\end{tikzpicture}%
}\\[0.2cm]
\raggedright\scriptsize \textbf{Fuente:} Elaboración propia a partir de García-Peñalvo (2024) y Artopoulos y Lliteras (2024).
\end{tcolorbox}
\end{figure}

\begin{cajabeige}
\textbf{CONCLUSIÓN DEL CASO:}\\
La IA generativa puede acelerar el trabajo, pero la responsabilidad de verificar, justificar y decidir sigue siendo estrictamente humana.
\end{cajabeige}

\newpage
\markboth{CAPÍTULO 1\quad NOTAS BIBLIOGRÁFICAS E HISTÓRICAS}{}
\subsection*{Notas bibliográficas e históricas}
\addcontentsline{toc}{subsection}{Notas bibliográficas e históricas}

\begin{itemize}
    \item \textbf{Sobre la definición y los conceptos (1.1):} Feuerriegel et al. (2024) ofrecen una introducción a la IA generativa en sus tres niveles, y García-Peñalvo y Vázquez-Ingelmo (2023) mapean su evolución y técnicas; ambos son un excelente punto de partida.
    \item \textbf{Sobre la historia (1.2):} Los textos originales de Turing (1950) y de la propuesta de Dartmouth (McCarthy et al., 2006) son breves y accesibles. Nasser-Eddine (2024) recorre las etapas y los inviernos de la IA. Las referencias clave de la historia reciente son Vaswani et al. (2017), Brown et al. (2020), Ho et al. (2020) y Ouyang et al. (2022).
    \item \textbf{Sobre el aprendizaje (1.3):} LeCun, Bengio y Hinton (2015) es la síntesis clásica del aprendizaje profundo y la retropropagación en \textit{Nature}.
    \item \textbf{Sobre la IA generativa y sus límites (1.4):} Bommasani et al. (2021) y Zhao et al. (2023) revisan los modelos fundacionales y de lenguaje; Goodfellow et al. (2014) y de la Torre (2023) tratan las GAN; Ji et al. (2023) y Bender et al. (2021) analizan alucinaciones y riesgos sociotécnicos; Sánchez y Carbajal (2023), García-Peñalvo (2024), García-Peñalvo et al. (2024) y Artopoulos y Lliteras (2024) abordan la educación superior.
\end{itemize}

\newpage
\markboth{CAPÍTULO 1\quad RESUMEN}{}
\subsection*{Resumen del Capítulo 1}
\addcontentsline{toc}{subsection}{Resumen del Capítulo 1}

La inteligencia artificial es un conjunto de métodos diseñados para realizar tareas propias de la cognición humana, evolucionando desde los primeros modelos lógicos hasta la actual era generativa. Este recorrido se sintetiza en seis puntos fundamentales:

\begin{resumenlist}
    \item[01] La IA busca construir sistemas que realicen tareas propias de la inteligencia humana; toda la IA actual en uso práctico es estrecha.
    \item[02] Su historia alterna avances e "inviernos": Turing y Dartmouth, el perceptrón, los sistemas expertos, el aprendizaje profundo y la era generativa.
    \item[03] Una máquina aprende ajustando los parámetros de un modelo para reducir un error medido sobre datos, mediante retropropagación y descenso del gradiente.
    \item[04] La IA generativa crea contenido nuevo; se apoya principalmente en Transformers, redes generativas antagónicas y modelos de difusión.
    \item[05] Los modelos de lenguaje predicen el siguiente token; su escala y la retroalimentación humana explican su utilidad práctica.
    \item[06] Persisten límites intrínsecos: alucinaciones, falta de comprensión real, sesgos, costos ambientales y riesgo de dependencia acrítica.
\end{resumenlist}

\begin{cajabeige}
\textbf{Síntesis:}\\
Definimos la IA como un conjunto de métodos para tareas que asociamos con la inteligencia humana, de tipo estrecho, y no como una mente ni una fuente infalible. Entre inviernos y veranos, despegó cuando coincidieron más datos, más cómputo y mejores algoritmos.
\end{cajabeige}

\newpage
\markboth{CAPÍTULO 1\quad CONCLUSIONES}{}
\subsection*{Conclusiones del Capítulo 1}
\addcontentsline{toc}{subsection}{Conclusiones del Capítulo 1}

\subsubsection*{Conclusión principal}
La IA generativa es la continuación de una historia larga, no una ruptura repentina. Las ideas de neuronas artificiales, aprendizaje a partir de datos y retropropagación existían desde hace décadas; lo que cambió fue la combinación de arquitecturas como el Transformer, grandes volúmenes de datos y cómputo masivo.

\subsubsection*{Aportes y límites}
Los modelos generativos amplían lo que una máquina puede hacer, pero no la vuelven infalible. Entender cómo aprenden y por qué fallan permite usarlos con criterio: verificar, citar y decidir con responsabilidad. La brecha entre el crecimiento técnico y los problemas aún no resueltos es, justamente, el tema que los siguientes capítulos permiten profundizar.

\subsubsection*{Preguntas abiertas}
\begin{enumerate}
    \item A medida que los modelos crecen y alucinan menos, \textquestiondown seguirá siendo suficiente que cada persona verifique lo que recibe, o harán falta mecanismos colectivos e institucionales para garantizar la confiabilidad?
    \item \textquestiondown Qué significará "crear" cuando una máquina pueda producir textos, cálculos e imágenes originales avanzadas, y cómo habría que evaluar el aprendizaje en ese escenario?
\end{enumerate}

\newpage
\markboth{CAPÍTULO 1\quad REFERENCIAS}{}
\subsection*{Referencias del Capítulo 1}
\addcontentsline{toc}{subsection}{Referencias del Capítulo 1}

\begin{itemize}[leftmargin=1.5em, label={}, itemsep=0.45ex, parsep=0pt]
    \item Artopoulos, A., \& Lliteras, A. (2024). Alfabetización crítica en IA. Recursos educativos para una pedagogía de la descajanegrización. \textit{Trayectorias Universitarias}, 10(19), e168. \url{https://doi.org/10.24215/24690090e168}
    \item Bender, E. M., Gebru, T., McMillan-Major, A., \& Shmitchell, S. (2021). On the dangers of stochastic parrots: Can language models be too big? \textit{Proceedings of the 2021 ACM FAccT}, 610--623. \url{https://doi.org/10.1145/3442188.3445922}
    \item Bommasani, R., Hudson, D. A., Adeli, E., Altman, R., et al. (2021). On the opportunities and risks of foundation models. \textit{arXiv:2108.07258}. \url{https://doi.org/10.48550/arXiv.2108.07258}
    \item Brown, T. B., Mann, B., Ryder, N., Subbiah, M., et al. (2020). Language models are few-shot learners. \textit{Advances in Neural Information Processing Systems (NeurIPS)}, 33, 1877--1901. \url{https://doi.org/10.48550/arXiv.2005.14165}
    \item de la Torre, J. (2023). Redes generativas adversarias (GAN): Fundamentos teóricos y aplicaciones. \textit{arXiv:2302.09346}. \url{https://doi.org/10.48550/arXiv.2302.09346}
    \item Feuerriegel, S., Hartmann, J., Janiesch, C., \& Zschech, P. (2024). Generative AI. \textit{Business \& Information Systems Engineering}, 66(1), 111--126. \url{https://doi.org/10.1007/s12599-023-00834-7}
    \item García-Peñalvo, F. J. (2024). Inteligencia artificial generativa y educación: Un análisis desde múltiples perspectivas. \textit{Education in the Knowledge Society}, 25, e31942. \url{https://doi.org/10.14201/eks.31942}
    \item García-Peñalvo, F. J., Llorens-Largo, F., \& Vidal, J. (2024). La nueva realidad de la educación ante los avances de la inteligencia artificial generativa. \textit{RIED}, 27(1), 9--39. \url{https://doi.org/10.5944/ried.27.1.37716}
    \item García-Peñalvo, F. J., \& Vázquez-Ingelmo, A. (2023). What do we mean by GenAI? \textit{IJIMAI}, 8(4), 7--16. \url{https://doi.org/10.9781/ijimai.2023.07.006}
    \item Goodfellow, I., Pouget-Abadie, J., Mirza, M., Xu, B., Warde-Farley, D., Ozair, S., Courville, A., \& Bengio, Y. (2014). Generative adversarial nets. \textit{Advances in Neural Information Processing Systems (NeurIPS)}, 27, 2672--2680. \url{https://doi.org/10.48550/arXiv.1406.2661}
    \item Ho, J., Jain, A., \& Abbeel, P. (2020). Denoising diffusion probabilistic models. \textit{Advances in Neural Information Processing Systems (NeurIPS)}, 33, 6840--6851. \url{https://doi.org/10.48550/arXiv.2006.11239}
    \item Ji, Z., Lee, N., Frieske, R., Yu, T., Su, D., Xu, Y., Ishii, E., Bang, Y. J., Madotto, A., \& Fung, P. (2023). Survey of hallucination in natural language generation. \textit{ACM Computing Surveys}, 55(12), 1--38. \url{https://doi.org/10.1145/3571730}
    \item LeCun, Y., Bengio, Y., \& Hinton, G. (2015). Deep learning. \textit{Nature}, 521(7553), 436--444. \url{https://doi.org/10.1038/nature14539}
    \item McCarthy, J., Minsky, M. L., Rochester, N., \& Shannon, C. E. (2006). A proposal for the Dartmouth Summer Research Project on Artificial Intelligence, August 31, 1955. \textit{AI Magazine}, 27(4), 12--14. \url{https://doi.org/10.1609/aimag.v27i4.1904}
    \item Nasser-Eddine López, F. (2024). Un repaso a la historia de la inteligencia artificial. \textit{Tabanque}, 36, 7--18. \url{https://doi.org/10.24197/trp.36.2024.7-18}
    \item Ouyang, L., Wu, J., Jiang, X., Almeida, D., et al. (2022). Training language models to follow instructions with human feedback. \textit{Advances in Neural Information Processing Systems (NeurIPS)}, 35, 27730--27744. \url{https://doi.org/10.48550/arXiv.2203.02155}
    \item Sánchez Mendiola, M., \& Carbajal Degante, E. (2023). La inteligencia artificial generativa y la educación universitaria. \textit{Perfiles Educativos}, 45(Especial), 70--86. \url{https://doi.org/10.22201/iisue.24486167e.2023.Especial.61692}
    \item Turing, A. M. (1950). Computing machinery and intelligence. \textit{Mind}, 59(236), 433--460. \url{https://doi.org/10.1093/mind/LIX.236.433}
    \item Vaswani, A., Shazeer, N., Parmar, N., Uszkoreit, J., Jones, L., Gomez, A. N., Kaiser, \L., \& Polosukhin, I. (2017). Attention is all you need. \textit{Advances in Neural Information Processing Systems (NeurIPS)}, 30, 5998--6008. \url{https://doi.org/10.48550/arXiv.1706.03762}
    \item Zhao, W. X., Zhou, K., Li, J., Tang, T., et al. (2023). A survey of large language models. \textit{arXiv:2303.18223}. \url{https://doi.org/10.48550/arXiv.2303.18223}
\end{itemize}
